\begin{abstract}
    This work investigates whether launch conditions and environmental parameters can be inferred from simulated three-dimensional projectile trajectories. A dataset of $200000$ trajectories is generated from a model incorporating gravity, quadratic drag, constant wind, and a simplified Magnus acceleration. Two regression approaches are compared: multilayer perceptrons operating on engineered trajectory descriptors and gated recurrent unit networks operating on resampled trajectory sequences. The models predict either the Cartesian initial velocity or a joint set of $11$ launch and environmental parameters. In addition to complete, noise-free trajectories, the sequence models are evaluated under partial observation and additive Gaussian position noise. Engineered features recover the restricted initial-velocity target with very high accuracy, whereas GRU models are more effective for the joint inverse problem. Supplying numerically derived velocity and acceleration channels improves parameter recovery on clean trajectories, and random-crop augmentation substantially improves predictions from short trajectory prefixes. However, these derivative channels strongly amplify measurement noise, causing marked degradation even for small positional perturbations. The results demonstrate both the potential of sequence models for physics-informed parameter inference and the importance of matching preprocessing and training augmentation to the expected observation conditions.
\end{abstract}
